% GENERATO da R/02_eiopa_rfr_bootstrap.R
\begin{table}[H]\centering\small
\caption{Metodo di Newton sull'equazione non lineare~\eqref{eq:boot-nl} per il gap $13\to15$ ($g=2$). Inizializzazione $f_0$ = forward annuale del tenor 13; criterio d'arresto $|f_k-f_{k-1}|\le 10^{-16}$. Il residuo $\varphi(f_k)$ crolla di diversi ordini di grandezza per iterazione: \`e la firma della convergenza quadratica.}
\label{tab:newton-esempio}\renewcommand{\arraystretch}{1.15}
\begin{tabular}{rccc}
\toprule
$k$ & $f_k$ (\%) & $\varphi(f_k)$ & $|f_k-f_{k-1}|$\\
\midrule
\rowcolor{rowtint}0 & 3.73842478 & -2.788e-03 & ---\\
1 & 3.58968675 & +5.977e-06 & 1.487e-03\\
\rowcolor{rowtint}2 & 3.59000428 & +2.734e-11 & 3.175e-06\\
3 & 3.59000428 & +0.000e+00 & 1.453e-11\\
\rowcolor{rowtint}4 & 3.59000428 & +0.000e+00 & 0.000e+00\\
\bottomrule
\end{tabular}
\end{table}
